% ============================================================
%  Trabajo Grupal 3 — LRD y Multifractales
%  Herramientas Matemáticas para el Manejo de la Información
%  Universidad Distrital Francisco José de Caldas — 2026-1
% ============================================================
\documentclass[12pt,a4paper]{article}

% ── Codificación y lenguaje ──────────────────────────────────
\usepackage[utf8]{inputenc}
\usepackage[T1]{fontenc}
\usepackage[spanish,shorthands=off]{babel}
\addto\captionsspanish{\renewcommand{\tablename}{Tabla}}

% ── Matemáticas ──────────────────────────────────────────────
\usepackage{amsmath,amssymb,amsthm}
\usepackage{mathtools}

% ── Geometría y tipografía ───────────────────────────────────
\usepackage{geometry}
\geometry{top=2.5cm, bottom=2.5cm, left=3cm, right=2.5cm, headheight=14pt}
\usepackage[protrusion=true,expansion=false]{microtype}
\usepackage{setspace}
\onehalfspacing

% ── Gráficos y flotantes ─────────────────────────────────────
\usepackage{graphicx}
\usepackage{float}
\usepackage{caption}
\usepackage{subcaption}

\usepackage{booktabs}
\usepackage{listings}
\usepackage{xcolor}
\usepackage{enumitem}
\usepackage{parskip}

% ── Listings (MATLAB) ───────────────────────────────────────
\lstdefinestyle{matlab}{
    language=Matlab,
    basicstyle=\ttfamily\small,
    keywordstyle=\color{blue}\bfseries,
    commentstyle=\color{green!60!black}\itshape,
    stringstyle=\color{red},
    numbers=left,
    numberstyle=\tiny\color{gray},
    stepnumber=1,
    numbersep=8pt,
    backgroundcolor=\color{gray!5},
    frame=single,
    rulecolor=\color{black!30},
    breaklines=true,
    breakatwhitespace=false,
    tabsize=4,
    captionpos=b,
    xleftmargin=1.5em,
    framexleftmargin=1.2em,
    showstringspaces=false
}
\lstset{style=matlab}

% ── TikZ ─────────────────────────────────────────────────────
\usepackage{tikz}
\usetikzlibrary{
    automata,
    positioning,
    arrows.meta,
    calc,
    shapes.geometric,
    shapes.misc,
    shadows.blur,
    backgrounds,
    decorations.pathreplacing
}
\usepackage{pgfplots}
\usepgfplotslibrary{groupplots,statistics}
\pgfplotsset{compat=1.18}

% ── Tablas ───────────────────────────────────────────────────
\usepackage{array}
\usepackage{longtable}
\usepackage{multirow}

% ── Hipervínculos ────────────────────────────────────────────
\usepackage[hidelinks]{hyperref}

% ── Entornos teóricos ────────────────────────────────────────
\theoremstyle{definition}
\newtheorem{definition}{Definición}[section]
\theoremstyle{plain}
\newtheorem{theorem}{Teorema}[section]
\newtheorem{proposition}{Proposición}[section]
\theoremstyle{remark}
\newtheorem{remark}{Observación}[section]

% ── Cabecera y pie ───────────────────────────────────────────
\usepackage{fancyhdr}
\pagestyle{fancy}
\fancyhf{}
\rhead{\small Trabajo Grupal 3}
\lhead{\small HMMI 2026-1}
\cfoot{\thepage}
\renewcommand{\headrulewidth}{0.4pt}

% ── Portada ──────────────────────────────────────────────────
\usepackage{titling}
\pretitle{%
    \begin{center}
    \includegraphics[width=0.40\textwidth]{assets/logo_circular_negro.pdf}\par\vspace{1em}
    \LARGE\bfseries
}
\posttitle{\par\end{center}\vspace{0.8em}}

% ============================================================
%  METADATOS
% ============================================================
\title{%
    \textbf{Trabajo Grupal 3}\\[4pt]
    \large LRD y Multifractales
    \vspace{2em}
}
\author{%
    \textit{Juan Felipe Rodríguez G.} \quad 20261595004 \\
    \textit{Carlos Stiven Mora Hoyos} \quad 20261595003 \\
    \vspace{3em}
    \textit{Sergio Santiago Duarte Rojas} \quad 20261595002 \\[6pt]
    \vspace{1em}
    Curso: Herramientas Matemáticas para el Manejo de la Información\\
    \vspace{1em}
    Universidad Distrital Francisco José de Caldas
    \vspace{1em}
}
\date{Bogotá, Colombia \\ \today}

% ============================================================
\begin{document}
\maketitle
\thispagestyle{fancy}

\clearpage
\tableofcontents
\clearpage

% ============================================================
\section{Actividad 1: Evaluación de estimadores con secuencias independientes}
% ============================================================

% Evaluar sesgo, consistencia, robustez y eficiencia del diagrama
% varianza-tiempo y del diagrama log-escala usando secuencias
% discretas independientes.

% ============================================================
\section{Actividad 2: Evaluación de estimadores con secuencias LRD}
% ============================================================

% Evaluar sesgo, consistencia y eficiencia del diagrama varianza-tiempo
% y del diagrama log-escala usando secuencias sintéticas con LRD
% generadas con wfbm y betamwm para varios valores de H (0.5 < H < 1).

% ============================================================
\section{Actividad 3: Estimación de H en la European ST-T Database}
% ============================================================

% Estimar el parámetro de Hurst del tiempo entre picos R de los 90
% registros. Graficar H vs. n con intervalos de confianza.
% Diagrama de caja y bigotes con boxplot.

% ============================================================
\section{Actividad 4: Herramienta de clasificación mono/multifractal}
% ============================================================

% Elaborar herramienta para estimar monofractalidad o multifractalidad
% de una secuencia con LRD. Validar con wfbm y betamwm.

% ============================================================
\section{Actividad 5: Clasificación fractal de la European ST-T Database}
% ============================================================

% Aplicar la herramienta de la Actividad 4 al tiempo entre picos R
% de los 90 registros. Analizar relación con las anotaciones de
% los especialistas.

% ============================================================
\section{Actividad 6: Escenario real con LRD}
% ============================================================

% Identificar un escenario real de la disciplina/profesión donde
% se presente LRD. Realizar análisis de fractalidad y concluir.

% ============================================================
\section*{Conclusiones}
% ============================================================

% TODO: Completar conclusiones.

% ============================================================
\begin{thebibliography}{99}

% TODO: Agregar referencias.

\end{thebibliography}

\end{document}
